\section{The character grading and the position of the Weil line}
\label{sec:characters}

This section proves \Cref{thm:character}, the obstruction announced in the
introduction and used in \Cref{prop:notpoly}. The group $\Kx$ acts on
$H^{2n}(A,\QQ)$ through the endomorphisms $\iota(\tau)$, this action grades
the cohomology, and the Weil line and the power $\eta^{n}$ of the polarisation
lie in different graded pieces. Once the grading is written down, the
obstruction follows in a few lines. The point that needs care is the
character through which $\Kx$ acts on the Weil line. It is
$\tau\mapsto\tau^{2n}$ together with its conjugate, and not a power of the
norm, although the Weil line is a $K$-line and multiplication by $\tau$ on a
$K$-line has determinant $\Nm(\tau)$ (\Cref{rem:failuremode}).

\subsection{The grading}

Throughout the section $(A,\iota,\eta)$ is of $(K,-1,n)$-Weil type,
$V=H^{1}(A,\QQ)$, and $V_{\CC}=V_{+}\oplus V_{-}$ is the splitting of
\Cref{not:eigen}.

\begin{proposition}[Character grading]\label{prop:grading}
For every $k\ge0$,
\begin{equation}\label{eq:grading}
  H^{k}(A,\CC)\;=\;\bigoplus_{r+s=k}
    \Bigl(\textstyle\bigwedge^{r}V_{+}\otimes\bigwedge^{s}V_{-}\Bigr),
\end{equation}
and $\iota(\tau)^{*}$ acts on the summand indexed by $(r,s)$ by the scalar
$\tau^{r}\bbar{\tau}^{\,s}$. The summands are the isotypic components for the
action of $\Kx$, and the characters $\tau\mapsto\tau^{r}\bbar{\tau}^{\,s}$ for
distinct pairs $(r,s)$ with $r+s=k$ are pairwise distinct.%
\footnote{Distinctness refers to characters of the torus
$\Res_{K/\QQ}\mathbb{G}_{m}$. At a fixed $\tau$ the values for $(r,s)$ and
$(r',s')$ agree if and only if $\tau/\bbar\tau$ is a root of unity of order
dividing $r-r'$, which can happen; see \Cref{rem:witness}.}
\end{proposition}

\begin{proof}
By \eqref{eq:wedge} we have $H^{k}(A,\CC)=\bigwedge^{k}V_{\CC}$, and the
exterior power of a direct sum decomposes as
$\bigwedge^{k}(V_{+}\oplus V_{-})=\bigoplus_{r+s=k}\bigwedge^{r}V_{+}\otimes\bigwedge^{s}V_{-}$,
which is \eqref{eq:grading}. By \Cref{not:eigen} the operator
$\iota(\tau)^{*}$ acts on $V_{+}$ by $\tau$ and on $V_{-}$ by $\bbar{\tau}$, so
on $\bigwedge^{r}V_{+}\otimes\bigwedge^{s}V_{-}$ it acts by
$\tau^{r}\bbar{\tau}^{\,s}$, since it is an algebra map for the wedge
product.

For the last statement, suppose $\tau^{r}\bbar{\tau}^{\,s}=\tau^{r'}\bbar{\tau}^{\,s'}$
for all $\tau\in\Kx$, with $r+s=r'+s'=k$. Dividing, $(\tau/\bbar{\tau})^{r-r'}=1$
for every $\tau\in\Kx$. If $r\ne r'$ this says that the image of the map
$c(\tau)=\tau/\bbar{\tau}$ lies in the group of $(r-r')$-th roots of unity of
$K$, which is finite because $K$ is an imaginary quadratic field and its group
of roots of unity has order $2$, $4$ or $6$. But $c$ has infinite image: two
elements have the same image exactly when their ratio lies in $\QQ^{\times}$,
and the elements $m+\sqrt{-d}$ for $m\in\ZZ$ are pairwise non-proportional over
$\QQ$. Hence $r=r'$ and so $s=s'$.
\end{proof}

\begin{lemma}[The polarisation sits in the middle]\label{lem:NSchar}
Let $(A,\iota,\eta)$ be of $(K,-1,n)$-Weil type. Then
\[
  \iota(\tau)^{*}\eta=\Nm_{K/\QQ}(\tau)\,\eta\qquad\text{for all }\tau\in\Kx,
\]
so $\eta$ lies in the summand $(r,s)=(1,1)$ of \eqref{eq:grading} and
$\eta^{n}$ lies in the summand $(r,s)=(n,n)$, on which $\Kx$ acts by $\Nm^{n}$.
\end{lemma}

\begin{proof}
Let $E$ be the alternating form on $H_{1}(A,\QQ)$ attached to $\eta$. Condition
(iii) of \Cref{def:weiltype} says that the Rosati involution $\dagger$ of
$\eta$ satisfies $\iota(\tau)^{\dagger}=\iota(\bbar{\tau})$, and by the
defining property of the Rosati involution,
\[
  E\bigl(\iota(\tau)x,\,y\bigr)=E\bigl(x,\,\iota(\tau)^{\dagger}y\bigr)
  =E\bigl(x,\,\iota(\bbar{\tau})y\bigr)
  \qquad\text{for all }x,y .
\]
Therefore
\begin{align*}
  \bigl(\iota(\tau)^{*}E\bigr)(x,y)
  &=E\bigl(\iota(\tau)x,\iota(\tau)y\bigr)
   =E\bigl(x,\iota(\bbar{\tau})\iota(\tau)y\bigr)\\
  &=E\bigl(x,\iota(\Nm(\tau))y\bigr)=\Nm(\tau)\,E(x,y),
\end{align*}
using $\bbar{\tau}\tau=\Nm(\tau)\in\QQ$ and the $\QQ$-linearity of $\iota$.
This is the displayed identity. By \Cref{prop:grading} the character $\Nm$ on
$H^{2}$ is $\tau\bbar{\tau}$, which is the summand $(r,s)=(1,1)$; and since
$\iota(\tau)^{*}$ is a ring map, $\iota(\tau)^{*}\eta^{n}=\Nm(\tau)^{n}\eta^{n}$,
placing $\eta^{n}$ in the summand $(n,n)$.
\end{proof}

\begin{lemma}[The Weil line sits at the ends]\label{lem:weilchar}
Let $(A,\iota,\eta)$ be of $(K,-1,n)$-Weil type. Then
\[
  \HW(A)\otimes_{\QQ}\CC
  =\Bigl(\textstyle\bigwedge^{2n}V_{+}\Bigr)\oplus
   \Bigl(\textstyle\bigwedge^{2n}V_{-}\Bigr),
\]
the summands $(r,s)=(2n,0)$ and $(0,2n)$ of \eqref{eq:grading} in degree $2n$,
on which $\Kx$ acts by $\tau\mapsto\tau^{2n}$ and $\tau\mapsto\bbar{\tau}^{\,2n}$.
\end{lemma}

\begin{proof}
This is \Cref{def:weilline} together with \Cref{prop:weilline}(i), which says
that $\HW(A)$ is the $\QQ$-structure on that pair of lines; the characters are
read off from \Cref{prop:grading} with $(r,s)=(2n,0)$ and $(0,2n)$.
\end{proof}

\begin{lemma}[The two characters differ]\label{lem:charsdiffer}
For every $n\ge1$ the characters $\tau\mapsto\tau^{2n}$ and
$\tau\mapsto\Nm_{K/\QQ}(\tau)^{n}=(\tau\bbar{\tau})^{n}$ of $\Kx$ are distinct,
as are $\tau\mapsto\bbar{\tau}^{\,2n}$ and $\Nm^{n}$.
\end{lemma}

\begin{proof}
The two characters agree if and only if $(\tau/\bbar{\tau})^{n}=1$ for every
$\tau\in\Kx$, since $\tau^{2n}=(\tau\bbar{\tau})^{n}$ is equivalent to
$(\tau/\bbar{\tau})^{n}=1$ for $\tau\ne0$.

Consider the map $c\colon\Kx\to\Kx$, $c(\tau)=\tau/\bbar{\tau}$. Two elements
have the same image if and only if their ratio is fixed by conjugation, that is
when it lies in $\QQ^{\times}$. The elements $m+\sqrt{-d}$ for $m\in\ZZ$ are
pairwise non-proportional over $\QQ$, so $c$ takes infinitely many values.

On the other hand $\{\zeta\in K:\zeta^{n}=1\}$ is contained in the group
$\mu(K)$ of roots of unity of $K$, and for an imaginary quadratic field
$\mu(K)$ is finite: it has order $4$ for $d=1$, order $6$ for $d=3$, and order
$2$ otherwise. An infinite set cannot be contained in a finite one, so some
$\tau\in\Kx$ has $(\tau/\bbar{\tau})^{n}\ne1$, and the characters differ. The
second pair is the image of the first under conjugation.
\end{proof}

\begin{remark}[Explicit witnesses]\label{rem:witness}
The proof does not exhibit a $\tau$ at which the two characters differ, and
the obvious candidate $\tau=1+\sqrt{-d}$ is not always one. For $d=1$ the
ratio $\tau/\bbar{\tau}$ is $\sqrt{-1}$, a primitive fourth root of unity, so
the two characters agree at this $\tau$ whenever $4\mid n$; for $d=3$ the
ratio is $(-1+\sqrt{-3})/2$, a primitive cube root of unity, and they agree
whenever $3\mid n$. The candidate $\tau=3+\sqrt{-d}$ fails as well, at $d=3$
with $6\mid n$, where the ratio is $(1+\sqrt{-3})/2$. Both failures of
$1+\sqrt{-d}$, with their exact exceptional sets, are certified in
\Cref{sec:lean}. The element $\tau=2+\sqrt{-d}$ is a witness for every $d$
and every $n$: its ratio is $\bigl(4-d+4\sqrt{-d}\bigr)/(4+d)$, which is not
$\pm1$, and for $d\notin\{1,3\}$ these are the only roots of unity in $K$,
while for $d=1$ and $d=3$ the ratios $(3+4\sqrt{-1})/5$ and
$(1+4\sqrt{-3})/7$ are not roots of unity either. The same holds for
$4+\sqrt{-d}$ and $3+2\sqrt{-d}$. The finiteness argument in the proof covers
all fields at once and needs no witness.
\end{remark}

\subsection{The obstruction}

\begin{proof}[Proof of \Cref{thm:character}]
The grading \eqref{eq:grading} and the action on its summands are
\Cref{prop:grading}. The placement of $\HW(A)$ is \Cref{lem:weilchar} and the
placement of $\eta^{n}$ is \Cref{lem:NSchar}. For the last assertion, write
$H^{2n}(A,\QQ)^{\Nm^{n}}$ for the space of $x$ with
$\iota(\tau)^{*}x=\Nm(\tau)^{n}x$ for all $\tau$. After $\otimes\CC$ this is
the isotypic component of the character $\Nm^{n}$, which by
\Cref{prop:grading} is the single summand $(r,s)=(n,n)$. By
\Cref{lem:charsdiffer} the characters $\tau^{2n}$ and $\bbar{\tau}^{\,2n}$ are
different from $\Nm^{n}$, so by \Cref{lem:weilchar} the space
$\HW(A)\otimes\CC$ meets the $(n,n)$ summand in zero. Intersecting with the
rational structure gives $\HW(A)\cap H^{2n}(A,\QQ)^{\Nm^{n}}=0$.
\end{proof}

\begin{corollary}[Equivariant classes]%
\label{cor:equivariant}
Let $x\in H^{2n}(A,\QQ)$ satisfy $\iota(\tau)^{*}x=\Nm(\tau)^{n}x$ for all
$\tau\in\Kx$. Then the component of $x$ along $\HW(A)$ in any $\Kx$-stable
complement is zero. In particular $x$ is a Weil class only if $x=0$.
\end{corollary}

\begin{proof}
Immediate from \Cref{thm:character}, since the isotypic decomposition is a
direct sum decomposition into $\Kx$-stable subspaces and $x$ lies in a single
summand that meets $\HW(A)$ in zero.
\end{proof}

\Cref{cor:equivariant} is the form of the obstruction that applies to
constructions. A construction that is natural in the $K$-action produces an
eigenclass, and the character of that eigenclass is determined by how the
input transforms. By \Cref{lem:NSchar,lem:charsdiffer}, the norm character,
which is the only one available to a rational eigenclass
(\Cref{rem:rightcharacter}), is the character already occupied by the
polarisation.

\begin{remark}[Rational eigenclasses]\label{rem:rightcharacter}
To land on the Weil line a construction must produce classes on which $\Kx$
acts through $\tau\mapsto\tau^{2n}$ or $\tau\mapsto\bbar{\tau}^{\,2n}$, and
neither character is $\QQ$-valued. This excludes rational eigenclasses. Let
$0\ne x\in H^{k}(A,\QQ)$ satisfy $\iota(\tau)^{*}x=\chi(\tau)\,x$ for all
$\tau\in\Kx$. Then $\chi(\tau)\in\QQ$, because $\iota(\tau)^{*}$ preserves
$H^{k}(A,\QQ)$. By \Cref{prop:grading}, $x$ lies in a single summand
$(r,s)$ and $\chi(\tau)=\tau^{r}\bbar{\tau}^{\,s}$; since $\chi$ is
$\QQ$-valued it equals its conjugate $\tau^{s}\bbar{\tau}^{\,r}$, and the
distinctness of the characters in \Cref{prop:grading} gives $r=s$. So every
rational $\Kx$-eigenclass in $H^{2n}(A,\QQ)$ is an eigenclass for $\Nm^{n}$,
and \Cref{cor:equivariant} applies to it. A construction that reaches the
Weil line must therefore produce a class that is not a $\Kx$-eigenclass. A
nonzero element of $\HW(A)$ has nonzero components in both summands
$(2n,0)$ and $(0,2n)$, since complex conjugation interchanges them, and its
$\Kx$-orbit spans the $\QQ$-plane $\HW(A)$. The Weil line is the one place in
the grading that no rational eigenvalue condition isolates.
\end{remark}

\begin{remark}[Determinant and character]\label{rem:failuremode}
The opposite conclusion can be reached by a plausible chain of statements.
The Weil line is $\bigwedge^{2n}_{K}V$, a free $K$-module of rank one; such a
module is two-dimensional over $\QQ$; and multiplication by $\tau$ on it is a
$\QQ$-linear map of determinant $\Nm(\tau)$. Each statement is true, but the
last concerns the $K$-module structure and not the action of $\iota(\tau)$.
By \Cref{lem:weilchar}, $\iota(\tau)$ acts on the Weil line as multiplication
by $\tau^{2n}$, with determinant $\Nm(\tau)^{2n}$, and a space on which $\Kx$
acts through multiplication by $\tau^{2n}$ is not a space on which it acts by
the scalar $\Nm(\tau)^{n}$. The first is a $K$-line and the second a
$\QQ$-eigenspace; they would agree only if $\tau^{n}=\bbar{\tau}^{\,n}$ for all
$\tau\in\Kx$, which \Cref{lem:charsdiffer} excludes. \Cref{fig:charactergrid}
draws the distinction.
\end{remark}

\begin{figure}[htbp]
\centering
  \includegraphics[width=\textwidth]{fig_charactergrid}
\caption{The $\Kx$-grading at $n=3$ on the grid of pairs $(r,s)$,
$0\le r,s\le 6$. The summands $\bigwedge^{r}V_{+}\otimes\bigwedge^{s}V_{-}$ of
$H^{6}(A,\CC)$ lie on the antidiagonal $r+s=6$ and are drawn at their
dimensions $\binom{6}{r}\binom{6}{s}$, which sum to $\binom{12}{6}=924$; on
the summand $(r,s)$ the group $\Kx$ acts by $\tau^{r}\bbar{\tau}^{\,s}$. The
two end cells $(6,0)$ and $(0,6)$, with characters $\tau^{6}$ and
$\bbar{\tau}^{\,6}$, span the Weil line; the middle cell $(3,3)$, with
character $\Nm(\tau)^{3}$, contains $\eta^{3}$. The outlined cells on the
diagonal $r=s$ are the summands on which $\Kx$ acts through a power of the
norm, and they carry the powers $\eta^{k}$. That the end cells and the middle
cell differ is the content of \Cref{thm:character}.}
\label{fig:charactergrid}
\end{figure}
